\par\Needspace{8\baselineskip}\bigskip
\pdfbookmark[3]{Darstellung, 16}{pallas-II-article16}
\begin{center}16.\end{center}
\noindent\hspace*{\parindent}Möchte man anstelle von \(T,V\) lieber \(R,S\) verwenden, so ist es für die Rechnung offenbar am zweckmäßigsten, die Orte des störenden wie des gestörten Planeten auf drei zueinander senkrechte Ebenen zu beziehen. Eine davon ist die Ebene der oskulierenden Bahn des gestörten Planeten, und eine der beiden anderen geht durch dessen Apsidenlinie. Es seien \(X',Y',Z'\) die Koordinaten des störenden Planeten und \(X,Y,0\) die des gestörten Planeten, und zwar
\SourceLeaf{466}
\[
\begin{aligned}
X'&=r'\cos PL',&\qquad X&=r\cos(v-\tilde\omega)\\
Y'&=r'\cos\Pi L',&Y&=r\sin(v-\tilde\omega)\\
Z'&=r'\cos ML',&&
\end{aligned}
\]
dann gilt
\[
\begin{gathered}
\rho=\BareRoot{((X'-X)^2+(Y'-Y)^2+Z'Z')}\\[4pt]
\begin{aligned}
\dfrac{R}{\mu}&=\dfrac{X'-X}{\rho^3}-\dfrac{X'}{r'^2}\\
\dfrac{S}{\mu}&=\dfrac{Y'-Y}{\rho^3}-\dfrac{Y'}{r'^3}\\
\dfrac{W}{\mu}&=\dfrac{Z'}{\rho^3}-\dfrac{Z'}{r'^3}\mathpunct{\cdot}
\end{aligned}
\end{gathered}
\]
\par Das zweckmäßigste Verfahren zur Berechnung dieser Koordinaten ist hinreichend bekannt, sodass wir hier darauf nicht näher eingehen müssen.
\par In allen vorangehenden Untersuchungen ist vorausgesetzt, dass die verschiedenen Bögen in Teilen des Radius ausgedrückt sind. In allen homogenen Formeln, in denen also die Bögen auf beiden Seiten dieselbe Dimension haben, kann man sie deshalb auch in Bogensekunden ausdrücken. In den inhomogenen Formeln muss man den Faktor 206265 auf der Seite einführen, auf der die Bögen eine Dimension weniger haben. Die Gleichungen 61, 62, 66, 68, 70, 74, 75, 76 bleiben daher unverändert. In den Gleichungen 63, 64, 65, 72, 73 muss man die rechte Seite durch 206265 teilen, wenn alle Bögen in Bogensekunden ausgedrückt werden. Nimmt man dagegen den Bogen \(n\) davon aus und drückt ihn weiterhin in Teilen des Radius aus, so ist die rechte Seite der Gleichungen 61, 62, 66 usw. mit 206265 zu multiplizieren, während die übrigen Gleichungen 63, 64, 65 usw. unverändert bleiben.
\par\medskip\begin{center}\rule{4em}{.4pt}\end{center}
